\chapter{仮説が間違っていた時の確信 $\Omega$ の棄却と前提を疑うダブルフィードバックループ学習}

\section{ベイズ更新の限界：シングルループにおける誤差修正の行き止まり}

ベイズ更新において、観測データ $D_1, D_2, \dots, D_m$ に基づいて事後オッズを累積計算していくプロセスは、事前に設定された「仮説空間 $\mathcal{H} = \{H_1, H_2, \dots, H_K\}$」の内部でどの仮説が最も妥当かを絞り込む作業である。
しかし、現場のプロファイリングや臨床介入において、どれほど精緻にエビデンスをサンプリングしても累積デシバンスコアが判定閾値（$+20\text{ dB}$）に到達しない、あるいは一度確定させた確信 $\Omega$ に反する強烈な陰性エビデンス（$-27\text{ dB}$）が突如として突きつけられる局面が存在する。

従来の制御理論や組織学習論において、クリス・アージリス（Chris Argyris）らはこれを\textbf{「シングルループ学習（Single-Loop Learning）」}と\textbf{「ダブルループ学習（Double-Loop Learning）」}の差異として定式化した。

\begin{itemize}
  \item \textbf{シングルループ（単一フィードバック）}：\\
  「既存の枠組み・前提（Governing Variables）」を一切疑わず、その枠組みの内部で行動やパラメータ（オッズ計算）の微修正のみを繰り返すフィードバック。
  \item \textbf{ダブルループ（二重フィードバック）}：\\
  誤差や矛盾が生じた際、変数そのものではなく、\textbf{「そもそも、その計算の基礎となっている前提、仮説空間の切り方、暗黙のルールそのもの」}を上位メタ階層から疑い、再定義するフィードバック。
\end{itemize}

サイコ・サイバネティクスおよび疑似ベイズ推論機がエゴ（DMN）の暴走に陥る最大の要因は、このシングルループの閉塞にある。
「相手は嘘をついているはずだ（仮説 $H_{\mathrm{lie}}$）」という固定された仮説空間の中でいくらベイズ更新を回しても、相手が潔白である場合、推論機はノイズの中から無理やり都合の良いシグナルを拾い集め、確信 $\Omega$ を捏造しようとする。
破綻を回避するためには、仮説が否定された瞬間に\textbf{「確信 $\Omega$ を冷徹に即時破棄（Purge）」}し、上位のダブルフィードバックループを起動させなければならない。

\section{オッズ $\le 1 : 100$（$-20\text{ dB}$）による確信 $\Omega$ の機械的棄却}

Intellectual Zero Personality における第一の規律は、オセロ・フィルターの判定基準に従った「冷徹な損切り（棄却）」である。

累積デシバンスコアが棄却閾値である \textbf{$-20\text{ dB}$（事後オッズ $1 : 100$ 以下、事後確率 $1\%$ 未満）} を割り込んだ場合、あるいは強力な陰性エビデンスによって事後オッズが急落した場合、推論機は直ちに以下のコマンドを実行する。

\begin{equation}
    \Omega_{\mathrm{current}} \xrightarrow[\text{Score} \le -20\text{ dB}]{\text{Hard Reset}} \emptyset \quad (\text{Null: 確信の完全蒸発})
\end{equation}

未熟な観察者やエゴに支配された人間は、「せっかくここまで立てた仮説を捨てたくない」「自分の直感が間違っていたと認めたくない」という保身バイアス（DMNの防衛）に囚われ、棄却閾値を超えても仮説に執着する。
しかし、本システムにおいて仮説とは「作業用の使い捨てツール」にすぎない。
$-20\text{ dB}$ に達した仮説は即座にメモリからパージされ、推論状態は完全な保留状態（$\mid$）へとリセットされる。

\section{ダブルフィードバックループの作動：メタ前提の解体と再編成}

仮説 $H$ が棄却されたとき、真に問われるべきは「では、対立仮説 $\bar{H}$ のどれが正しいのか？」という単純な二者択一ではない。
システムは直ちに第2のループ（上位フィードバック）へと制御権を移行させる。

\begin{equation}
\begin{array}{rcccl}
\text{\textbf{[第2ループ]}} & \multicolumn{3}{c}{\boldsymbol{\text{根本前提・コンテクストの再定義} \ (\mathcal{M} \to \mathcal{M}')}} & \nwarrow \\
& \downarrow & & & \quad \Big\uparrow \ \text{\textbf{仮説空間の破綻}} \\
\text{\textbf{[第1ループ]}} & \text{仮説空間 } \mathcal{H} & \xrightarrow[\text{Bayes Update}]{D_1, D_2, \dots} & \text{事後オッズ計算} & \longrightarrow \text{スコア} \le -20\text{ dB}
\end{array}
\end{equation}

このとき、CEN（中央実行ネットワーク）が精査すべき「疑うべき前提（メタ・パラメータ）」は以下の3点である。

\begin{enumerate}[label=\textbf{(\arabic*)}]
  \item \textbf{仮説空間の排他性・網羅性の疑義（カテゴリーエラーの検出）}：\\
  「嘘をついているか（$H_1$）、真実を述べているか（$H_2$）」という二項対立そのものが間違っていなかったか？
  「相手自身が記憶を解離させており、客観的には嘘だが主観的には真実だと思い込んでいる（第3の仮説 $H_3$）」という可能性はないか？
  既存の枠組みから漏れていた新たな次元を導入し、仮説空間 $\mathcal{H}$ そのものを拡張・再生成する。
  
  \item \textbf{尤度モデル $P(D \mid H)$ の前提条件（ベースライン）の疑義}：\\
  オセロ・フィルターの尤度比を計算する前提となった「ベースライン（平熱状態）」の測定が狂っていなかったか？
  例えば、相手が重度の睡眠不足、慢性疼痛、あるいはカフェインの過剰摂取状態にあった場合、通常の人間なら「やや稀（$P=0.1$）」である微小硬直や発汗が、その個体にとっては「ベースライン（$P=0.5$）」であった可能性がある。
  生体ハードウェアの物理状態を再キャリブレーションし、尤度パラメータを再設定する。

  \item \textbf{コンテクスト（文脈）設定のズレ}：\\
  術者が投げかけたプロービング（微小入力）が、術者の意図した意味とは全く異なる文脈として相手の疑似ベイズ推論機にデコードされていなかったか？
  言葉の定義や文脈のフレーム（枠組み）をメタ・モデルを用いて再質問し、通信プロトコルのハンドシェイクをやり直す。
\end{enumerate}

\section{工学的総括：確信への執着をゼロにする「動的再構成マシン」}

シングルループに閉じこもる人間は、現実が自分の仮説と矛盾したとき、「現実を歪めて自分の確信 $\Omega$ を守ろう」とする（これが妄想、オセロの誤謬、カルト的盲信の正体である）。

対照的に、ダブルフィードバックループを実装した Intellectual Zero Personality は、現実と仮説の矛盾を「前提を再定義するための極上の情報（メタ・エビデンス）」として歓迎する。
\begin{center}
\textbf{「仮説の棄却とは敗北ではなく、仮説空間の次元を引き上げるためのトリガーである」}
\end{center}
確信 $\Omega$ をミリ秒単位で形成し、不要になればナノ秒単位で棄却して上位前提を書き換える。
この動的なダブルループが確立されて初めて、生体コンピュータは環境の変化と他者の深層に対して、完全にバイアスフリーな無限の追従性を獲得するのである。

